% !TeX root = ../main.tex
% 本LaTeX模板的一般使用说明
\chapter{说明}

Again，这是北航论文\LaTeX{}模板（\CTeX{}-Based）\BUAAThesis{}。

本\LaTeX{}模板为北航研究生学位论文模板，适用于理工类博士、硕士。本\LaTeX{}模板参考自《北京航空航天大学研究生撰写学位论文的规定》，根据2021年11月修订版调整，具体要求请参见各自的《北京航空航天大学研究生撰写学位论文的规定》，最终成文格式需参考学院要求及打印方意见。


%-----------------------------
\section{宏包使用}

请将以下文件与此LaTeX文件放在同一目录中：

\begin{tabular}{ll}
 \verb|def/buaa.cls |                 & $\triangleright$ LaTeX宏模板文件 \\
 \verb|def/GBT7714-2015.bst|      & $\triangleright$ 国标参考文献BibTeX样式文件2015 \\
 \verb|def/GBT7714-2015-NoWarning.bst|      & $\triangleright$ 国标参考文献BibTeX样式文件2015且不提示缺失信息 \\
 \verb|def/simsun.ttc|            & $\triangleright$ 宋体 \\
 \verb|def/logo-buaa.eps|         & $\triangleright$ 论文封皮北航字样 \\
 \verb|def/head-master.eps|       & $\triangleright$ 论文封皮学术硕士学位论文标题\\
 \verb|def/head-professional.eps| & $\triangleright$ 论文封皮专业硕士学位论文标题\\
 \verb|def/head-doctor.eps|       & $\triangleright$ 论文封皮学术博士学位论文标题\\
 \verb|def/head-prodoctor.eps|    & $\triangleright$ 论文封皮专业博士学位论文标题\\
 \verb|tex/*.tex|                 & $\triangleright$ 本模板样例中的独立章节\\
\end{tabular}\\

通过 \verb|\documentclass[| \verb|<thesis>,| \verb|<permission>,| \verb|<printtype>,| \verb|<ostype>,| \verb|<ctexbookoptions>| \verb|]{buaa}|载入宏包：
\begin{itemize}[leftmargin=3cm]
  \item[{\tt thesis} $\triangleright$]  论文类型(thesis)，可选值：\\
    a) 学术硕士论文（\verb|master|）[缺省值]\\
    b) 专业硕士论文（\verb|professional|）\\
    c) 学术博士论文（\verb|doctor|）\\
    d) 专业博士论文（\verb|prodoctor|）
  \item[{\tt permission} $\triangleright$] 密级(permission)，可选值： \\
    a) 公开（\verb|public|）[缺省值]\\
    b) 内部（\verb|privacy|）\\
    c) 秘密（\verb|secret|=\verb|secret3|）\\
    --- c.1) 秘密3年（\verb|secret3|）\\
    --- c.2) 秘密5年（\verb|secret5|）\\
    --- c.3) 秘密10年（\verb|secret10|）\\
    --- c.4) 秘密永久（\verb|secret*|）\\
    d) 机密（\verb|classified|=\verb|classified5|）\\
    --- d.1) 机密3年（\verb|classified3|）\\
    --- d.2) 机密5年（\verb|classified5|）\\
    --- d.3) 机密10年（\verb|classified10|）\\
    --- d.4) 机密永久（\verb|classified*|）\\
    e) 绝密（\verb|topsecret|=\verb|topsecret10|）\\
    --- e.1) 绝密3年（\verb|topsecret3|）\\
    --- e.2) 绝密5年（\verb|topsecret5|）\\
    --- e.3) 绝密10年（\verb|topsecret10|）\\
    --- e.4) 绝密永久（\verb|topsecret*|）
  \item[{\tt printtype} $\triangleright$] 打印设置(printtype)，可选值： \\
    a) 单面打印（\verb|oneside|）[缺省值]\\
    b) 双面打印（\verb|twoside|）
  \item[{\tt ostype} $\triangleright$] 系统类型(printtype)，可选值： \\
    a) Windows（\verb|win|）[缺省值]\\
    b) Linux（\verb|linux|）\\
    c) Mac（\verb|mac|）
  \item[{\tt ctexbookoptions} $\triangleright$] {\tt ctexbook}文档类支持的其他选项： \\
    使用{\tt ctexbookoptions}选项传递{\tt ctexbook}文档类支持的其他选项。
    例如，使用{\tt fontset=founder}选项启用方正字体以避免生僻字乱码的问题\footnote{需要系统安装方正字体。}。
\end{itemize}


\setlength{\hangindent}{4em}
模板已内嵌LaTeX工具包：\\
{\tt ifthen}，{\tt etoolbox}，{\tt titletoc}，{\tt remreset}，
{\tt geometry}，{\tt fancyhdr}，{\tt setspace}，{\tt float}，
{\tt graphicx}，{\tt subfigure}，{\tt epstopdf}，{\tt array}，{\tt enumitem}，
{\tt booktabs}，{\tt longtable}，{\tt multirow}，{\tt caption}，
{\tt listings}，{\tt algorithm2e}，{\tt amsmath}，{\tt amsthm}，
{\tt hyperref}，{\tt pifont}，{\tt color}，{\tt soul}；\\
For Windows: {\tt times}, {\tt newtxmath}；\\
For Linux: {\tt newtxtext}, {\tt newtxmath}；\\
For Mac: {\tt times}, {\tt fontspec}。


模板已内嵌宏：\verb|\highlight{text}|（黄色高亮）。

请统一使用UTF-8编码。



%-----------------------------
\section{选项设置}

\begin{itemize}[leftmargin=3cm]
  \item[{\tt  $\backslash$refcolor} $\triangleright$]  开启/关闭引用编号颜色，包括参考文献，公式，图，表，算法等\\
  \texttt{on}：开启 [默认]\\
  \texttt{off}：关闭
  \item[{\tt $\backslash$beginright} $\triangleright$]  摘要和正文从右侧页开始\\
  \texttt{on}：开启 [默认]\\
  \texttt{off}：关闭
  \item[{\tt $\backslash$emptypageword} $\triangleright$]  空白页留字
  \item[{\tt $\backslash$Listfigtab} $\triangleright$]  是否使用图标清单目录\\
  \texttt{on}：开启 [默认]\\
  \texttt{off}：关闭
\end{itemize}


%-----------------------------
\section{章节撰写}
本模板支持以下标题级别标题级别：

\begin{tabular}{ll}
  \verb|\chapter{章}|              & $\triangleright$ 第一章 \\
  \verb|\chapter*{无章号章}|       & $\triangleright$ 无章号章 \\
  \verb|\chaptera{无章号有目录章}| & $\triangleright$ 无章号有目录章 \\
  \verb|\summary|                  & $\triangleright$ 总结\\
  \verb|\appendix|                 & $\triangleright$ 附录\\
  \verb|\achievement|              & $\triangleright$ 攻读学位期间取得的成果\\
  \verb|\acknowledgments|          & $\triangleright$ 致谢\\
  \verb|\biography|                & $\triangleright$ 作者简介\\
  \verb|\section{节}|              & $\triangleright$ 1.1 节\\
  \verb|\subsection{条}|           & $\triangleright$ 1.1.1 条\\
  \verb|\subsubsection{A}|         & $\triangleright$ 1.1.1.1 A\\
  \verb|\paragraph{a}|             & $\triangleright$ 1.1.1.1.1 a\\
  \verb|\subparagraph{a)}|         & $\triangleright$ 1.1.1.1.1.1 a)\\
\end{tabular}

%-----------------------------
\section{注意事项}
\begin{itemize}
  \item[$\triangleright$] \textit{中文斜体}将转换为楷体；
  \item[$\triangleright$] \textbf{中文粗体}在Windows（From WeiQM）和Mac（From CaiBW）下转换为黑体，Linux下正常（From QiaoJF）；
  \item[$\triangleright$] \verb|\label{<text>}|中不能使用中文；
  \item[$\triangleright$] 浮动体与正文之间的距离是弹性的，根据内容调整，不太好控制；
  \item[$\triangleright$] 命令符与汉字之间请注意加空格以避免undefined错误（pdfLaTeX下好像一般不存在这个问题，主要在XeLaTeX编译环境下发生）；
\end{itemize}

%-----------------------------
\section{ToDo}
\begin{itemize}
  \item[$\triangleright$] 数学环境的行间隔；
  \item[$\triangleright$] 参考文献的行间隔；
\end{itemize}

%-----------------------------
\section{意见及问题反馈}

\indent 请作答该问卷：https://www.wjx.cn/vm/hxvWNe5.aspx\\

\begin{figure}[!h]
	\centering
	\includegraphics[width=.5\textwidth]{pic/question_survey.jpg}
	\caption{问题反馈问卷二维码}
	\label{fig:survey_ques}
\end{figure}